%%%%%%%%%%%%%%%%%%%%%%%%%%%%%%%%%%%%%%%%%%%%%%%%%%%%%%%%%%%%%%%%%%%%%%%%%%%%%%%
%                                                                             %
% BlackArch Linux Guide                                                       %
% Translated by: SOVALO                                                       %
%                                                                             %
%%%%%%%%%%%%%%%%%%%%%%%%%%%%%%%%%%%%%%%%%%%%%%%%%%%%%%%%%%%%%%%%%%%%%%%%%%%%%%%

\documentclass[a4paper, oneside, 11pt]{book}

%%% INCLUDES %%%
\renewcommand{\familydefault}{\sfdefault}

\usepackage{array}
\usepackage{color}
\usepackage{enumerate}
\usepackage{fancyhdr}
\usepackage{fancyvrb}
\usepackage{geometry}
\usepackage{graphicx}
\usepackage{html}
\usepackage{hyperref}
\usepackage{ifpdf}
\usepackage{listings}
\usepackage{pstricks}
\usepackage{supertabular}
\usepackage{tocloft}

\usepackage{fontspec}
\setmainfont{CMU Serif}[Ligatures=TeX]
\setsansfont{CMU Sans Serif}[Ligatures=TeX]
\usepackage[ukrainian]{babel}
\usepackage[utf8]{inputenc}
\usepackage[T1]{fontenc}

\renewcommand{\contentsname}{Зміст}
\renewcommand{\chaptername}{Розділ}

%%% LAYOUT %%%
\setlength{\parindent}{0em}
\setlength{\parskip}{1.5ex plus0.5ex minus0.5ex}
\geometry{left=2.5cm, textwidth=16cm, top=3cm, textheight=25cm, bottom=3cm}
\widowpenalty=2000
\clubpenalty=1000
\frenchspacing
\sloppy
\pagecolor[HTML]{FFFFFF}
\color[HTML]{333333}
\setcounter{tocdepth}{10}
\setcounter{secnumdepth}{10}

\hypersetup{
  pdftitle={BlackArch Linux, Посібник користувача BlackArch Linux},
  pdfsubject={BlackArch Linux, Посібник користувача BlackArch Linux},
  pdfauthor={BlackArch Linux, BlackArch Linux},
  pdfkeywords={BlackArch Linux, Penetration Testing, Security, Hacking, Linux},
  pdfcenterwindow=true,
  colorlinks=true,
  breaklinks=true,
  linkcolor=blue,
  menucolor=blue,
  urlcolor=red
}

% syntax highlighting
% all options should be set here document wide
\lstset{
backgroundcolor=\color[HTML]{EEEEEE},
frame=single,
basicstyle=\footnotesize\ttfamily,
float,
deletekeywords={return},
otherkeywords={mkdir, curl, sudo, sha1sum, grep, cut, sort, wget, makepkg,
pacman, blackman, chmod},
keywordstyle=\color{orange},
commentstyle=\color{blue},
stringstyle=\color{red},
language=bash,
showspaces=false,
showtabs=false,
tabsize=2
}

%%% HEADER / FOOTER %%%
\setlength{\headheight}{33pt}
\setlength{\headsep}{33pt}
\lhead{{\includegraphics[width=1cm,height=1cm]{images/logo.png}}}
\rhead{Посібник користувача BlackArch Linux}

%%% CUSTOM MACROS %%%
% for html links
\ifpdf\else
\def\href#1#2{\htmladdnormallink{#2}{#1}}
\fi

%------------------%
%   TITLE PAGE     %
%------------------%
\begin{document}
\pagestyle{empty}
\begin{center}
\begin{figure}[htbp]
\centering
\vspace{0.5cm}
\includegraphics[width=8cm]{images/logo.png}
\label{fig:logo}
\end{figure}
\vspace{0.5cm}
\Huge{\textbf{The BlackArch Linux Guide}}\\
\vspace{1cm}
\Large{\color{blue}https://www.blackarch.org/}\\
\vspace{0.5cm}
\end{center}
\newpage
\tableofcontents
\newpage
\pagestyle{fancy}

%------------------%
%  Chapter 1       %
%------------------%

\chapter{Вступ}

\section{Огляд}
Посібник BlackArch Linux розділений на кілька частин:
\begin{itemize}
\item Вступ --- містить загальний огляд, знайомство та додаткову корисну інформацію про проєкт.
\item Посібник користувача --- усе, що необхідно знати звичайному користувачеві для ефективного використання BlackArch.
\item Посібник розробника --- як розпочати розробку та зробити свій внесок у BlackArch.
\item Посібник з інструментів --- детальна інформація про утиліти на практичних прикладах використання (В розробці).
\end{itemize}

\section{Що таке BlackArch Linux?}
BlackArch --- це повноцінний дистрибутив Linux для тестувальників на проникнення та дослідників безпеки. Він заснований на \href{https://www.archlinux.org/}{ArchLinux}, і користувачі можуть встановлювати компоненти BlackArch окремо або цілими групами.

Набір інструментів розповсюджується як \href{https://wiki.archlinux.org/index.php/Unofficial\_User\_Repositories}{неофіційний користувацький репозиторій} Arch Linux, тому ви можете встановити BlackArch поверх уже наявної системи Arch Linux. Пакети можна встановлювати індивідуально або за категоріями.

Репозиторій постійно розширюється і наразі налічує понад \href{https://www.blackarch.org/tools.html}{2600} інструментів. Усі утиліти ретельно тестуються перед додаванням до кодової бази для підтримання належної якості репозиторію.

\section{Історія BlackArch Linux}
Незабаром...

\section{Підтримувані платформи}
Незабаром...

\section{Долучитися до проєкту}
Ви можете зв'язатися з командою BlackArch за допомогою таких каналів:

Вебсайт: \url{https://www.blackarch.org/}

Email: \href{mailto:team@blackarch.org}{team@blackarch.org}

IRC: \url{irc://irc.freenode.net/blackarch}

Twitter: \url{https://twitter.com/blackarchlinux}

GitHub: \url{https://github.com/Blackarch/}

Discord: \url{https://discord.com/invite/xMHt8dW}

%------------------%
%  Chapter 2       %
%------------------%

\chapter{Посібник користувача}

\section{Встановлення}
У наступних розділах описано, як налаштувати репозиторій BlackArch та встановити пакети. BlackArch підтримує обидва варіанти: встановлення з репозиторію за допомогою бінарних пакетів, а також їхнє компілювання та встановлення з вихідного коду.

BlackArch повністю сумісний зі звичайним встановленням Arch Linux, оскільки виступає в ролі неофіційного користувацького репозиторію. Якщо замість цього вам потрібен повноцінний ISO-образ, див. розділ \href{https://www.blackarch.org/downloads.html#iso}{Live ISO}.

\subsection{Встановлення поверх наявної системи ArchLinux}
Запустіть скрипт \href{https://blackarch.org/strap.sh}{strap.sh} з правами суперкористувача (root) та дотримуйтесь інструкцій. Див. приклад нижче:
\begin{lstlisting}
   curl -O https://blackarch.org/strap.sh
   sha1sum strap.sh # має збігатися з: 5ea40d49ecd14c2e024deecf90605426db97ea0c
   sudo chmod +x strap.sh
   sudo ./strap.sh
\end{lstlisting}

Тепер завантажте свіжу копію головного списку пакетів та виконайте синхронізацію системи:
\begin{lstlisting}
  sudo pacman -Syyu
\end{lstlisting}

\subsection{Встановлення пакетів}
Тепер ви можете встановлювати інструменти з репозиторію BlackArch.
\begin{enumerate}
\item Щоб переглянути всі доступні інструменти, виконайте:
\begin{lstlisting}
  pacman -Sgg | grep blackarch | cut -d' ' -f2 | sort -u
\end{lstlisting}

\item Щоб встановити взагалі всі інструменти, виконайте:
\begin{lstlisting}
  pacman -S blackarch
\end{lstlisting}

\item Щоб встановити конкретну категорію інструментів, виконайте:
\begin{lstlisting}
  pacman -S blackarch-<category>
\end{lstlisting}

\item Щоб переглянути доступні категорії BlackArch, виконайте:
\begin{lstlisting}
  pacman -Sg | grep blackarch
\end{lstlisting}
\end{enumerate}

\subsection{Встановлення пакетів із вихідного коду}
Альтернативний метод дозволяє збирати пакети BlackArch безпосередньо з вихідного коду. Ви можете знайти файли PKGBUILD на сторінці \href{https://github.com/BlackArch/blackarch/tree/master/packages}{GitHub}. Для збирання всього репозиторію можна скористатися утилітою \href{https://github.com/BlackArch/blackman}{Blackman}.
\begin{itemize}
\item По-перше, вам потрібно встановити Blackman. Якщо у вашій системі вже налаштовано репозиторій пакетів BlackArch, ви можете встановити його через pacman:
\begin{lstlisting}
  pacman -S blackman
\end{lstlisting}

\item Також ви можете зібрати та встановити Blackman з вихідного коду самостійно:
\begin{lstlisting}
  mkdir blackman
  cd blackman
  wget https://raw.github.com/BlackArch/blackarch/master/packages/blackman/PKGBUILD
  # Переконайтеся, що файл PKGBUILD не був зловмисно модифікований.
  makepkg -s
\end{lstlisting}

\item Або ви можете встановити Blackman через будь-який AUR-помічник:
\begin{lstlisting}
  <ваш AUR-помічник> -S blackman
\end{lstlisting}
\end{itemize}

\subsection{Базове використання Blackman}
Blackman дуже простий у використанні, хоча його прапорці відрізняються від того, до чого ви зазвичай звикли в pacman. Основні команди наведено нижче:
\begin{itemize}
\item Завантажити, скомпілювати та встановити конкретний пакет:
\begin{lstlisting}
  sudo blackman -i package
\end{lstlisting}

\item Завантажити, скомпілювати та встановити цілу категорію:
\begin{lstlisting}
  sudo blackman -g group
\end{lstlisting}

\item Завантажити, скомпілювати та встановити взагалі всі інструменти BlackArch:
\begin{lstlisting}
  sudo blackman -a
\end{lstlisting}

\item Вивести список категорій BlackArch:
\begin{lstlisting}
  blackman -l
\end{lstlisting}

\item Вивести список інструментів у конкретній категорії:
\begin{lstlisting}
  blackman -p category
\end{lstlisting}
\end{itemize}

\subsection{Встановлення з образів Live-ISO, Netinstall-ISO або ArchLinux}
Ви можете встановити BlackArch Linux, використовуючи один із наших офіційних образів Live чи Netinstall ISO. Див. \url{https://www.blackarch.org/download.html#iso}. Після завантаження з ISO необхідно виконати такі кроки:
\begin{itemize}
\item Встановіть пакет інсталятора:
\begin{lstlisting}
  sudo pacman -S blackarch-installer
\end{lstlisting}

\item Запустіть процес встановлення:
\begin{lstlisting}
  sudo blackarch-install
\end{lstlisting}
\end{itemize}

%------------------%
%  Chapter 3       %
%------------------%

\chapter{Посібник розробника}

\section{Система збирання та репозиторіїв Arch}
Файли PKGBUILD --- це скрипти збирання. Кожен із них вказує утиліті makepkg(1), як саме створити пакет. Файли PKGBUILD пишуться мовою Bash.

Для отримання детальнішої інформації ознайомтеся з такими ресурсами:
\begin{itemize}
\item \href{https://wiki.archlinux.org/index.php/Creating_Packages}{Arch Wiki: Creating Packages}
\item \href{https://wiki.archlinux.org/index.php/Makepkg}{Arch Wiki: makepkg}
\item \href{https://wiki.archlinux.org/index.php/PKGBUILD}{Arch Wiki: PKGBUILD}
\item \href{https://wiki.archlinux.org/index.php/Arch_Packaging_Standards}{Arch Wiki: Arch Packaging Standards}
\end{itemize}

\section{Стандарти PKGBUILD у BlackArch}
Задля простоти, наші файли PKGBUILD дуже схожі на стандарти AUR, за винятком кількох невеликих відмінностей, описаних нижче. Кожен пакет має обов'язково належати щонайменше до групи `blackarch`; багато пакетів можуть входити до складу кількох груп одночасно.

\subsection{Групи}
Щоб дозволити користувачам швидко та легко завантажувати певний спектр інструментів, пакеты були розділені на тематичні групи. Наявність груп дозволяє користувачеві просто виконати команду \verb|pacman -S <назва групи>|, щоб отримати набір необхідного софту.

\subsubsection{blackarch}
Це основна група, до якої входять абсолютно всі пакети дистрибутиву. Вона дозволяє користувачам легко встановити кожен наявний інструмент.

Що має бути тут: усе підряд.

\subsubsection{blackarch-anti-forensic}
Пакети, що використовуються для протидії криміналістичному аналізу та цифровій експертизі (forensics), включаючи шифрування, стеганографію та будь-які засоби, що змінюють атрибути файлів чи системи з метою приховування інформації.

Приклади: luks, TrueCrypt, Timestomp, dd, ropeadope, secure-delete

\subsubsection{blackarch-automation}
Пакети, що використовуються для автоматизації різноманітних робочих процесів (workflows).

Приклади: blueranger, tiger, wiffy

\subsubsection{blackarch-backdoor}
Пакети, які використовують вразливості або встановлюють бекдори (приховані входи) на вже скомпрометованих системах.

Приклади: backdoor-factory, rrs, weevely

\subsubsection{blackarch-binary}
Пакети, призначені для аналізу, модифікації чи обробки бінарних файлів у будь-якому вигляді.

Приклади: binwally, packerid

\subsubsection{blackarch-bluetooth}
Пакети, які використовуються для експлуатації вразливостей у пристроях та протоколах стандарту Bluetooth (802.15.1).

Приклади: ubertooth, tbear, redfang

\subsubsection{blackarch-code-audit}
Пакети для проведення аудиту наявного вихідного коду з метою пошуку потенційних вразливостей.

Приклади: flawfinder, pscan

\subsubsection{blackarch-cracker}
Пакети, призначені для зламу криптографічних функцій, підбору паролів та розшифрування хешів.

Приклади: hashcat, john, crunch

\subsubsection{blackarch-crypto}
Пакети, що працюють із криптографією, за винятком безпосереднього зламу чи підбору.

Приклади: ciphertest, xortool, sbd

\subsubsection{blackarch-database}
Пакети, пов'язані з експлуатацією та тестуванням баз даних на будь-якому рівні.

Приклади: metacoretex, blindsql

\subsubsection{blackarch-debugger}
Пакети (відладчики), які дозволяють користувачеві відстежувати та аналізувати дії певної програми в режимі реального часу.

Приклади: radare2, shellnoob

\subsubsection{blackarch-decompiler}
Пакети (декомпілятори), які використовуються для спроби відновлення вихідного коду вже скомпільованих програм.

Приклади: flasm, jd-gui

\subsubsection{blackarch-defensive}
Пакети, які використовуються для захисту системи користувача від шкідливого ПЗ та кібератак з боку інших користувачів.

Приклади: arpon, chkrootkit, sniffjoke

\subsubsection{blackarch-disassembler}
Пакети (дизассемблери), що перетворюють машинний код програми на читабельний текст мовою асемблера.

Приклади: inguma, radare2, smali

\subsubsection{blackarch-dos}
Пакети, що використовуються для проведення DoS/DDoS атак, метою яких є доведення системи чи сервісу до відмови в обслуговуванні.

Приклади: 42zip, nkiller2

\subsubsection{blackarch-drone}
Пакети, призначені для перехоплення керування чи взаємодії з фізично сконструйованими дронами та БПЛА.

Приклади: meshdeck, skyjack

\subsubsection{blackarch-exploitation}
Пакети, що безпосередньо експлуатують відомі чи знайдені вразливості в інших програмах, операційних системах або мережевих службах.

Приклади: armitage, metasploit, zarp

\subsubsection{blackarch-fingerprint}
Пакети, що використовуються для зняття "відбитків" (fingerprinting) --- ідентифікації версій ОС, ПЗ, мережевих пристроїв та конфігурацій.

Приклади: dns-map, p0f, httprint

\subsubsection{blackarch-firmware}
Пакети, призначені для пошуку та експлуатації вразливостей безпосередньо у прошивках (мікропрограмах) пристроїв.

Приклади: Наразі немає, будуть додані найближчим часом.

\subsubsection{blackarch-forensic}
Пакети для проведення цифрової криміналістичної експертизи, пошуку та відновлення даних на фізичних накопичувачах або в оперативній пам'яті.

Приклади: aesfix, nfex, wyd

\subsubsection{blackarch-fuzzer}
Пакети, що використовують принцип фаззинг-тестування (fuzzing), який полягає в надсиланні додатку некоректних, неочікуваних або випадкових вхідних даних для перевірки стабільності та пошуку багів.

Приклади: msf, mdk3, wfuzz

\subsubsection{blackarch-hardware}
Пакети, які використовуються для взаємодії, керування чи експлуатації фізичного апаратного забезпечення.

Приклади: arduino, smali

\subsubsection{blackarch-honeypot}
Пакети, які діють як приманки (honeypots) --- симулюють вразливі служби чи системи для приваблення зловмисників та вивчення їхньої тактики.

Приклади: artillery, bluepot, wifi-honey

\subsubsection{blackarch-keylogger}
Пакети, призначені для запису та збереження натискань клавіш на цільовій системі.

Приклади: Наразі немає, будуть додані найближчим часом.

\subsubsection{blackarch-malware}
Пакети, що використовуються для виявлення або аналізу будь-яких типів шкідливого програмного забезпечення.

Приклади: malwaredetect, peepdf, yara

\subsubsection{blackarch-misc}
Різноманітні пакети загального призначення, які не вписуються чітко в жодну іншу категорію.

Приклади: oh-my-zsh-git, winexe, stompy

\subsubsection{blackarch-mobile}
Пакети для роботи з мобільними платформами та смартфонами.

Приклади: android-sdk-platform-tools, android-udev-rules

\subsubsection{blackarch-networking}
Пакети, пов'язані з аналізом, налаштуванням та роботою з IP-мережами на будь-якому рівні.

Приклади: Майже будь-які мережеві утиліти.

\subsubsection{blackarch-nfc}
Пакети для експлуатації та роботи з технологією бездротового зв'язку малого радіуса дії NFC (near-field communications).

Приклади: nfcutils

\subsubsection{blackarch-packer}
Пакети, що працюють з пакувальниками (packers) або пов'язані з ними.
\textit{Пакувальники --- це програми, які маскують чи стискають шкідливе ПЗ всередині інших виконуваних файлів для обходу антивірусів.}

Приклади: packerid

\subsubsection{blackarch-proxy}
Пакети, що функціонують як проксі-сервери, тобто перенаправляють мережевий трафик через проміжні вузли.

Приклади: burpsuite, ratproxy, sslnuke

\subsubsection{blackarch-recon}
Пакети, які активно шукають вразливості в реальних умовах (активна розвідка). Здебільшого це загальна група для схожих за призначенням інструментів інформаційного збору.

Приклади: canri, dnsrecon, netmask

\subsubsection{blackarch-reversing}
Загальна зонтична група для будь-яких декомпіляторів, дизассемблерів або інструментів зворотного інжинірингу (reverse engineering).

Приклади: capstone, radare2, zerowine

\subsubsection{blackarch-scanner}
Пакети, призначені для сканування цільових мереж та систем з метою виявлення відкритих портів та вразливостей.

Приклади: scanssh, tiger, zmap

\subsubsection{blackarch-sniffer}
Пакети, які використовуються для перехоплення та глибокого аналізу мережевого трафіку (сніфінгу).

Приклади: hexinject, pytactle, xspy

\subsubsection{blackarch-social}
Пакети, орієнтовані на проведення атак, збір інформації або тестування безпеки соціальних мереж та платформ.

Приклади: jigsaw, websploit

\subsubsection{blackarch-spoof}
Пакети, призначені для підміни даних (спуфінгу) --- наприклад, IP-адрес чи MAC-адрес, щоб приховати справжнього автора запитів від жертви.

Приклади: arpoison, lans, netcommander

\subsubsection{blackarch-threat-model}
Пакети, що використовуються для фіксації, аналізу та моделювання загроз безпеки у конкретно визначених сценаріях.

Приклади: magictree

\subsubsection{blackarch-tunnel}
Пакети, що застосовуються для інкапсуляції (тунелювання) мережевого трафіку всередині інших протоколів.

Приклади: ctunnel, iodine, ptunnel

\subsubsection{blackarch-unpacker}
Пакети, які використовуються для вилучення раніше упакованого чи замаскованого шкідливого коду з виконуваних файлів.

Приклади: js-beautify

\subsubsection{blackarch-voip}
Пакети для тестування та експлуатації програм і протоколів IP-телефонії (VoIP).

Приклади: iaxflood, rtp-flood, teardown

\subsubsection{blackarch-webapp}
Пакети, призначені для сканування та тестування безпеки вебдодатків і сайтів, що дивляться в інтернет.

Приклади: metoscan, whatweb, zaproxy

\subsubsection{blackarch-windows}
Ця група містить утиліти, створені спеціально під ОС Windows, але які стабільно працюють у Linux через прошарок Wine.

Приклади: 3proxy-win32, pwdump, winexe

\subsubsection{blackarch-wireless}
Пакети для аудиту та експлуатації бездротових мереж на будь-якому рівні (Wi-Fi тощо).

Приклади: airpwn, mdk3, wiffy

\sectionsection{Структура репозиторію}
Головний Git-репозиторій BlackArch розташований за адресою: \\
\url{https://github.com/BlackArch/blackarch}. \\
\\Також існує кілька додаткових репозиторіїв екосистеми: \\
\url{https://github.com/BlackArch}.

У головному репозиторії є три ключові каталоги:
\begin{itemize}
\item docs --- документація та посібники.
\item packages --- файли збирання PKGBUILD.
\item scripts --- корисні допоміжні скрипти автоматизації.
\end{itemize}

\subsection{Скрипти}
Нижче наведено короткий опис утиліт, розташованих у каталозі \verb|scripts/|:
\begin{itemize}
  \item baaur --- інструмент для автоматичного надсилання пакетів до AUR (у розробці).
  \item babuild --- скрипт для безпосереднього збирання пакета.
  \item bachroot --- керування оточенням chroot для безпечного тестування пакетів.
  \item baclean --- очищення застарілих файлів .pkg.tar.xz із локального сховища пакетів.
  \item baconflict --- заміна для застарілого скрипта \verb|scripts/conflicts| (у розробці).
  \item bad-files --- пошук некоректних чи пошкоджених файлів у зібраних пакетах.
  \item balock --- отримання або зняття блокування (lock) з репозиторію пакетів.
  \item banotify --- надсилання сповіщень на IRC-канал про додавання нових пакетів.
  \item barelease --- публікація готових пакетів безпосередньо у публічний репозиторій.
  \item baright --- відображення інформації про авторські права та ліцензії BlackArch.
  \item basign --- цифровий підпис зібраних пакетів.
  \item basign-key --- підпис ключів безпеки репозиторію.
  \item blackman --- працює подібно до pacman, але збирає софт напряму з джерел через git (не плутати з оригінальною утилітою Blackman від nrz).
  \item check-groups --- перевірка правильності призначення груп для пакетів.
  \item checkpkgs --- тестування зібраних пакетів на наявність поширених помилок.
  \item conflicts --- перевірка файлової структури на наявність можливих конфліктів.
  \item dbmod --- модифікація та оновлення бази даних пакетів репозиторію.
  \item depth-list --- створення списку пакетів, відсортованого за глибиною їхніх залежностей.
  \item deptree --- побудова дерева залежностей виключно для пакетів, що надаються BlackArch.
  \item get-blackarch-deps --- виведення списку залежностей для конкретного пакета.
  \item get-official --- отримання офіційних пакетів для поточного релізу.
  \item list-loose-packages --- пошук пакетів, які не належать до жодної групи та не мають залежностей від інших компонентів.
  \item list-needed --- відображення списку відсутніх у системі залежностей.
  \item list-removed --- виведення пакетів, які досі є в бінарному репозиторії, але вже видалені з git.
  \item list-tools --- генерація повного списку підтримуваних інструментів.
  \item outdated --- пошук пакетів, версії яких у бінарному репозиторії застаріли порівняно з актуальним станом у git.
  \item pkgmod --- внесення модифікацій у конфігурацію збирання пакетів.
  \item pkgrel --- автоматичне збільшення індексу релізу (pkgrel) у файлі пакета.
  \item prep --- автоматичне очищення стилю коду у файлі PKGBUILD та базовий пошук синтаксичних помилок.
  \item sitesync --- синхронізація локальної копії бінарного репозиторію з віддаленим сервером дзеркал.
  \item size-hunt --- пошук пакетів, що займають критично багато дискового простору.
  \item source-backup --- резервне копіювання вихідних файлів та архівів.
\end{itemize}

\section{Внесення вкладу в проєкт (Contribution)}
У цьому розділі описано процес додавання власних напрацювань до BlackArch Linux. Ми приймаємо будь-які pull requests --- від виправлення дрібних одруків до додавання абсолютно нових інструментів. Якщо у вас виникли запитання чи пропозиції, ви завжди можете зв'язатися з нашою командою.
\\\\
Будь-хто може стати контриб'ютором. Будь-яка допомога є важливою та цінною для проєкту.

\subsection{Обов'язкові керівництва}
Будь ласка, уважно вивчіть ці інструкції перед тим, як надсилати свої перші зміни:
\begin{itemize}
\item \href{https://wiki.archlinux.org/index.php/Arch\_Packaging\_Standards)}{Arch Packaging Standards}
\item \href{https://wiki.archlinux.org/index.php/Creating\_Packages}{Creating Packages}
\item \href{https://wiki.archlinux.org/index.php/PKGBUILD}{PKGBUILD}
\item \href{https://wiki.archlinux.org/index.php/Makepkg}{Makepkg}
\end{itemize}

\subsection{Покроковий алгоритм дій}
Щоб запропонувати свої зміни до кодової бази BlackArch Linux, виконайте такі кроки:
\begin{enumerate}
\item Зробіть власний форк офіційного репозиторію: \url{https://github.com/BlackArch/blackarch}
\item Модифікуйте необхідні файли (файли PKGBUILD, патчі `.patch` тощо).
\item Закомітьте внесені зміни у свою гілку.
\item Надішліть (пушніть) зміни у свій віддалений форк на GitHub.
\item Створіть офіційний запит на злиття коду через механізм **Pull Request**.
\end{enumerate}

\subsection{Практичний приклад}
Наступний приклад демонструє процес додавання нового пакета до проєкту BlackArch. Ми скористаємося утилітою \href{https://wiki.archlinux.org/index.php/yaourt}{yaourt} (також підійде pacaur або будь-який інший AUR-helper) для завантаження наявного файлу PKGBUILD для утиліти \textbf{nfsshell} з репозиторію \href{https://aur.archlinux.org/}{AUR} з його подальшою адаптацією під наші вимоги.

\subsubsection{Отримання файлу PKGBUILD}
Завантажуємо файли пакета з AUR за допомогою інструменту yaourt:
\begin{lstlisting}
  user@blackarchlinux $ yaourt -G nfsshell
  ==> Download nfsshell sources
  x LICENSE
  x PKGBUILD
  x gcc.patch
  user@blackarchlinux $ cd nfsshell/
\end{lstlisting}

\subsubsection{Очищення стилю та форматування}
Очищаємо структуру та форматування файлу \textit{PKGBUILD} за допомогою нашого внутрішнього скрипта `prep`:
\begin{lstlisting}
  user@blackarchlinux nfsshell $ ./blackarch/scripts/prep PKGBUILD
  cleaning 'PKGBUILD'...
  expanding tabs...
  removing vim modeline...
  removing id comment...
  removing contributor and maintainer comments...
  squeeze extra blank lines...
  removing '|| return'...
  removing leading blank line...
  removing $pkgname...
  removing trailing whitespace...
\end{lstlisting}

\subsubsection{Редагування та налаштування}
Тепер відкриваємо файл \textit{PKGBUILD} у текстовому редакторі для внесення необхідних правок та категорій:
\begin{lstlisting}
  user@blackarchlinux nfsshell $ vi PKGBUILD
\end{lstlisting}

\subsubsection{Збирання пакета}
Запускаємо локальне збирання пакета для перевірки його працездатності:
\begin{lstlisting}
user@blackarchlinux nfsshell $ makepkg -sf
==> Making package: nfsshell 19980519-1 (Mon Dec  2 17:23:51 CET 2013)
==> Checking runtime dependencies...
==> Checking buildtime dependencies...
==> Retrieving sources...
-> Downloading nfsshell.tar.gz...
% Total    % Received % Xferd  Average Speed   Time    Time     Time
CurrentDload  Upload   Total   Spent    Left  Speed100 29213  100 29213    0
0  48150      0 --:--:-- --:--:-- --:--:-- 48206
-> Found gcc.patch
-> Found LICENSE
...
<сюди виводиться лог компіляції та збирання пакета>
...
==> Leaving fakeroot environment.
==> Finished making: nfsshell 19980519-1 (Mon Dec  2 17:23:53 CET 2013)
\end{lstlisting}

\subsubsection{Встановлення та тестування}
Встановлюємо свіжозібраний пакет у систему та тестуємо його працездатність:
\begin{lstlisting}
  user@blackarchlinux nfsshell $ pacman -U nfsshell-19980519-1-x86_64.pkg.tar.xz
  user@blackarchlinux nfsshell $ nfsshell # перевіряємо роботу програми
\end{lstlisting}

\subsubsection{Додавання, коміт та відправка змін}
Переносимо перевірену папку інструменту до каталогу пакетів репозиторію, створюємо коміт та пушимо зміни:
\begin{lstlisting}
user@blackarchlinux nfsshell $ cd /blackarchlinux/packages
user@blackarchlinux ~/blackarchlinux/packages $ mv ~/nfsshell .
user@blackarchlinux ~/blackarchlinux/packages $ git commit -am nfsshell && git push
\end{lstlisting}

\subsubsection{Створення Pull Request}
Перейдіть у браузер та відкрийте сторінку свого форку на GitHub, щоб створити Pull Request:
\begin{lstlisting}
  firefox https://github.com/<ваш_юзернейм>/blackarchlinux
\end{lstlisting}

\subsubsection{Налаштування віддаленого репозиторію (Upstream)}
Якщо ви плануєте працювати з проєктом довгостроково, гарною практикою буде додати оригінальний репозиторій як `upstream`, щоб синхронізувати свій локальний код:
\begin{lstlisting}
  user@blackarchlinux ~/blackarchlinux $ git remote -v
  origin <url вашого форку> (fetch)
  origin <url вашого форку> (push)
  user@blackarchlinux ~/blackarchlinux $ git remote add upstream https://github.com/blackarch/blackarch
  user@blackarchlinux ~/blackarchlinux $ git remote -v
  origin <url вашого форку> (fetch)
  origin <url вашого форку> (push)
  upstream https://github.com/blackarch/blackarch (fetch)
  upstream https://github.com/blackarch/blackarch (push)
\end{lstlisting}

За замовчуванням утиліта git буде відправляти зміни прямо в `origin` (ваш форк). Це правильна поведінка, адже у вас немає прямих прав на пуш в основну кодову базу проєкту без схвалення Pull Request.

Якщо ж ви з часом отримаєте статус офіційного розробника з правами запису, для зручності та надійності краще використовувати SSH-адресу вигляду \textit{git@github.com:blackarch/blackarch.git}.

\subsection{Основні вимоги до оформлення коду}
\begin{enumerate}
\item Не додавайте коментарі на кшталт \textbf{Maintainer} або \textbf{Contributor} безпосередньо всередину файлів \textit{PKGBUILD}. Імена всіх авторів та контриб'юторів зазначаються окремо в спеціальному розділі AUTHORS загального гайду BlackArch.
\item Задля підтримання однорідності стилю коду в усьому репозиторії, будь ласка, використовуйте відступи розміром у **два пробіли** (two-space indentation).
\end{enumerate}

\subsection{Загальні поради}
Використовуйте утиліту \href{http://wiki.archlinux.org/index.php/Namcap}{namcap}, вона чудово допомагає автоматично сканувати готові пакети на наявність прихованих помилок або некоректних ліцензій.

%------------------%
%  Chapter 4       %
%------------------%

\chapter{Посібник з інструментів}
Незабаром...

\section{Скоро буде доступно}
Незабаром...

%%% APPENDIX %%%
\appendix
\include{latex/appendix-uk}

\end{document}

%%% EOF %%%
